\section{Projective ball quotients and a Frobenius identity}\label{sec:arithmetic}

The theta construction has supplied complex one-form classes on compact
ball quotients.  To use the surface criterion, we must identify their
rational sources: for the sign function \(v\), the relevant class must
lie in the complex span of the images of \(U(v)_1\) under rational Hodge
maps from \(U(v)\).  This section and Section~\ref{sec:cm-sources}
prove the following extraction theorem.  We first recall its input
data from the theta construction.

Retain the notation of Section~\ref{sec:theta}.  Thus \(E\subset\Qbar\)
is Galois CM, \(F=E^+\), \(G_E=\Gal(E/\Q)\), and \(c\in G_E\) is
complex conjugation.  We assume \(s=[F:\Q]\geq2\).  The Hermitian
\(E/F\)-space \(V\) has dimension three, signature \((2,1)\) at
\(\mu=1\), and is positive definite at the other real places.
Write \(G=\SU(V)\), let \(D\) be its ball of negative lines, and retain
the fixed \(\mathcal O_E\)-lattice \(\Lambda\subset V\).  Fix an odd
function \(v:G_E\to\{1,-1\}\) with \(v(1)=1\), and put
\[
d(\tau)=v(\tau^{-1}).
\]
For one multiplicity line \(P\) with these compact-place signs, write
\(\Theta(\phi)\) for its theta one-form and abbreviate its finite
special-unitary action to \(\omega_f=\omega_{P,f}\).  We work with a
finite collection of inputs at a common principal level.  The auxiliary
prime will be chosen only after that level is fixed.

\Needspace{11\baselineskip}
\begin{theorem}[Degree-one CM extraction]\label{thm:cm-extraction}
Let \(v:G_E\to\{1,-1\}\) be odd with \(v(1)=1\), and use the
multiplicity-line theta construction prescribed by
\(d(\tau)=v(\tau^{-1})\). Given a finite collection of its finite
adelic Schwartz inputs, there is a common sufficiently deep principal
ball quotient \(S_0=\Gamma\backslash D\) on which every associated
theta one-form class belongs to
\begin{equation}\label{eq:cm-extraction-span}
\operatorname{span}_{\C}\left\{
f\bigl(U(v)_1\bigr):
f\in\Hom_{\mathrm{Hdg}}\bigl(U(v),H^1(S_0,\Q)\bigr)
\right\}.
\end{equation}
The maps in this formula are rational Hodge maps, and their images
are viewed after extension to \(\C\). The conclusion persists after
passage to a deeper principal level.
\end{theorem}

The quotient \(S_0\) will be a smooth connected projective surface.
Corollary~\ref{cor:finite-family-extraction} gives the version needed
for the surface criterion: finitely many families with different sign
functions satisfy these span conclusions on one projective principal
cover of any prescribed torsion-free principal base quotient.  The reference
moduli construction depends only on \(V\), so the same geometry
serves every multiplicity line.

We begin by realizing the principal quotients as components of a full
projective moduli space of polarized abelian varieties with \(E\)-action.
Keeping all components makes the isogeny correspondences into algebraic
Hecke operators on the product of their Albanese varieties.  Dense
ordinary loci then let us prove an integral Frobenius polynomial for
these operators.  Section~\ref{sec:cm-sources} combines that polynomial
with the active theta parameter to determine the Hodge types of the
generalized Hecke spaces containing the theta classes.  Scalar descent
then proves the extraction theorem.

\subsection{A projective moduli space for the reference three-space}\label{subsec:reference-moduli}

The reference abelian varieties use \(V\), not \(V\otimes_E P\).
For a negative line \(L\in D\), prescribe on the complexified homology
\[
\mathcal F_\mu=L,\qquad
\mathcal F_c=\operatorname{ann}(L)\subset V_\mu^*,
\]
where the Hermitian form identifies the conjugate summand with the
complex-linear dual.  At each other conjugate pair choose, once and
for all, the whole summand on one side and zero on the other.  These
annihilator prescriptions together form a Lagrangian.  An imaginary element
\(\delta_{\mathrm{ref}}\in E^\times\), whose signs are chosen by weak
approximation in \(F\), makes them positive for
\[
\psi_{\mathrm{ref}}(x,y)
 =\Tr_{E/\Q}\bigl(\delta_{\mathrm{ref}}h_V(x,y)\bigr).
\]
Multiplying \(\delta_{\mathrm{ref}}\) by a positive integer makes the
form integral on \(\Lambda\).

Here the homology filtration is the kernel of the Lie quotient.  If
\(J\) is the complex structure on real homology, then
\(z=u+iJu\) lies in that filtration and
\[
\psi_{\mathrm{ref}}(z,\overline z)/i
   =-2\psi_{\mathrm{ref}}(u,Ju).
\]
Thus \(-\psi_{\mathrm{ref}}\) is the Riemann form in the positivity
convention of Section~\ref{sec:theta}.  The corresponding abelian
varieties have dimension \(3s\), and their Lie ranks are
\begin{equation}\label{eq:reference-lie-ranks}
r_\mu=2,\qquad r_c=1,\qquad
\{r_\tau,r_{c\tau}\}=\{0,3\}
\quad\bigl(\tau\notin\{1,c\}\bigr).
\end{equation}
Fix their polarization type.

Choose \(N\) sufficiently divisible to give a fine abelian moduli
problem, a torsion-free principal lattice level, and a level fixing
the chosen finite theta inputs.  We will also require the level to
kill all possible unitary determinant roots of unity.  A marking
means a full ordered basis of \(N\)-torsion; we allow every pairing
matrix for that basis.  In particular, transporting it through an
isogeny of degree prime to \(N\) is allowed even when the isogeny
rescales the polarization pairing.

We retain all components so that the full isogeny correspondences can
be defined without assuming that the reference component is preserved.
Over a number field \(K_0\subset\Qbar\) containing \(E\), consider the
reduced scheme \(S\) of tuples
\[
(B,\lambda,\iota_E,\eta_N).
\]
Here \(B\) is an abelian variety of dimension \(3s\), \(\lambda\) is
a polarization of the fixed type,
\(\iota_E:\mathcal O_E\to\End(B)\) is a unital action whose Rosati
adjoint is \(c\), \(\eta_N\) is the full marking, and each Lie
summand has the particular rank fixed above (with values displayed
in \eqref{eq:reference-lie-ranks}).  We retain all connected
components of this tuple space.

\begin{lemma}\label{lem:reference-moduli}
The scheme \(S\) is of finite type, smooth, and projective over \(K_0\).
Its complex connected components are compact ball quotients.  For
sufficiently divisible \(N\), the component arising from \(\Lambda\)
is the quotient of \(D\) by its principal special-unitary level.
\end{lemma}

\begin{proof}
The Lie-rank condition need not identify a single rational Hermitian
isomorphism class.  We therefore construct the whole finite-type
action locus, identify its complex components as ball quotients, and
then prove projectivity and identify the component of \(\Lambda\).

\emph{The finite-type action locus.}
The fine-moduli theorem for polarized abelian varieties gives a
quasi-projective scheme with universal polarization and full level at
sufficiently large invertible level
\cite[Chapter~7, Definitions~7.1--7.2, Proposition~7.6, and
Theorem~7.9]{MumfordFogarty1982}.
In that theorem the polarization degree is fixed and the level is an
ordered torsion basis, with no prescribed pairing matrix.  After
inverting the polarization degree, the polarization kernel is finite
étale, so its geometric elementary-divisor type is locally constant.
We may therefore take the open-and-closed part of the required type.
If necessary, first apply the theorem to a sufficiently large divisor
of \(N\) prime to the polarization degree and then add the remaining
full \(N\)-basis as a finite étale level cover.  These constructions
work over a ring of integers with finitely many primes inverted, as
well as in characteristic zero.

We next impose the action by a finite-type construction.  Choose
finitely many ring generators \(b\) of \(\mathcal O_E\) and
parameterize their graphs in relative Hilbert schemes of the square
of the universal abelian scheme.  The degrees of the graphs in
question are bounded.  Indeed, put \(\alpha=1+\bar b b\).  It is a
nonzero totally positive element of \(F\).  Choose a positive integer
\(m\) with \(m\alpha^{-1}\in\mathcal O_E\).  In every such action
\(\alpha\) is an isogeny, and
\[
\deg(\alpha)\deg(m\alpha^{-1})=m^{6s}.
\]
The second factor is a positive integer, so \(\deg(\alpha)\) is
bounded independently of the point.

To translate this into a Hilbert bound, the normalized Poincaré bundle
on the universal abelian scheme times its dual, pulled back by
\((\id,\lambda)\), is a relatively ample bundle inducing \(2\lambda\).
The tensor product of its two pullbacks to the square restricts on the graph of
\(b\) to a bundle inducing \(2\lambda(1+\bar b b)\).  Its polarization
degree is bounded by the preceding identity.  Riemann--Roch on an
abelian variety of dimension \(3s\) gives the Hilbert polynomial
\(\chi(\mathcal L)n^{3s}\), with
\(\chi(\mathcal L)^2=\deg(\phi_{\mathcal L})\).  Only finitely many
polynomials can therefore occur.  On their graph loci impose the
homomorphism, ring, Rosati, and Lie-rank identities.  This gives a
finite-type action locus.

Over the characteristic-zero polarized base this locus is unramified.
One form of rigidity suffices: an infinitesimal difference of two
homomorphisms between fixed abelian schemes would give a homomorphism
from a proper connected abelian fiber to an additive tangent group,
and must vanish.  The locus is also proper over that base.  Given a
discrete valuation ring mapping to the polarized base and a lift of
its generic point to the action locus, each
generator homomorphism extends uniquely between the fixed abelian
schemes by the Néron mapping property
\cite[\S1.2, Definition~1 and Proposition~8]{BoschLutkebohmertRaynaud1990}.
Its graph remains a graph of the same Hilbert polynomial.  The ring
and Rosati identities extend by uniqueness.  Since the base contains
\(E\), the action algebra splits into idempotents in characteristic
zero, and their Lie summands have locally constant ranks.  Thus the
extended point lies in the same action locus.  The valuative criterion
gives properness.  A proper finite-type unramified morphism is
quasi-finite and hence finite, so the action locus is finite over the
polarized base.

\emph{The complex components.}
We now identify the reduced complex structure of this algebraic locus.
In a marked period chart,
the integral matrices of the action are locally constant.  Such a
matrix acts on the complex torus exactly when it preserves the
homology filtration; the resulting holomorphic endomorphism of a
polarized complex torus is algebraic.  Hence the reduced local locus
is precisely the locus of positive filtrations stable under the fixed
matrices.  Trace duality and the Rosati identity express the
alternating form as an imaginary trace multiple of an
\(E/F\)-Hermitian form on the rational homology, of \(E\)-dimension
three.  Isotropy makes the two filtrations in a conjugate pair
annihilators.  Positivity and \eqref{eq:reference-lie-ranks} leave one
negative line at the active place and constant data at every other
place.  The local locus is therefore a smooth complex two-ball.

This description gives whole ball quotients, not only local ball
charts.  Fix a marked polarized \(E\)-lattice occurring in \(S\).
Every positive filtration in its compatible ball gives a complex
torus with the fixed integral Riemann form, and hence an abelian
variety; the preserving integral matrices give its algebraic
\(\mathcal O_E\)-action.  Conversely every tuple with these lattice
data gives such a filtration.  Its connected component is thus the
quotient of the whole ball by the marked lattice stabilizer.

\emph{Projectivity and the principal level.}
For each such rational Hermitian space, a vector isotropic over \(E\)
would remain isotropic at a definite real place.  Since \(s\geq2\),
there is no such vector, and its special-unitary group is
\(F\)-anisotropic.  The arithmetic lattice and compactness theorems
give compact ball quotients
\cite{BorelHarishChandra1962}; see
\cite[Theorem~3.3 and Proposition~3.5]{MilneShimura2017}.
The level removes stabilizers.  There are finitely many components
because \(S\) is of finite type.  The finite action locus is
quasi-projective, and its reduced complex space is smooth with compact
components.  It is therefore proper, and a proper quasi-projective
scheme is projective.  Smoothness and properness descend from the
complex base change to \(K_0\).

Finally the initial lattice \(\Lambda\) is among these data.  Even
if it is not free, its determinant module is a fractional
\(\mathcal O_E\)-ideal.  The determinant of a unitary lattice
automorphism is consequently a unit of \(\mathcal O_E\) of norm one.
All its complex conjugates have modulus one, so it is a root of unity.
A sufficiently divisible principal level kills every nontrivial such
determinant.  Its stabilizer is then precisely the principal
special-unitary lattice level, as claimed.
\end{proof}

\subsection{Good models and ordinary points on every component}\label{subsec:good-ordinary-models}

Enlarge \(K_0\) to a number field \(K\subset\Qbar\) over which all
connected components \(S_j\) are geometrically connected and have
\(K\)-points.  Index the original ball component by \(S_0\).
Use those points as origins for the Albanese maps
\[
a_j:S_j\longrightarrow A_j=\Alb(S_j),\qquad A_S=\prod_j A_j.
\]
We regard each \(A_j\) as a factor of \(A_S\), with zero in the other
factors.  The Albanese pullbacks give an isomorphism of rational Hodge
structures
\begin{equation}\label{eq:albanese-h1}
H^1(A_S,\Q)\simeq H^1(S,\Q).
\end{equation}

Because the Frobenius identity will be checked on ordinary points and
then on Albanese differences, we need density on every component and
generation of each reduced Albanese factor by those differences.

\begin{lemma}\label{lem:good-ordinary-primes}
After one finite extension of \(K\) and removal of finitely many
primes, the \(S_j\) have smooth proper models with geometrically
connected fibers, the \(A_j\) extend to abelian schemes, and the
\(a_j\) extend to the models.  The model of \(S\) lies in the full
tuple parameter scheme, retaining the \(\mathcal O_E\)-action and its
graph coordinates.  Differences of Albanese images generate every
geometric fiber of each \(A_j\).

At every remaining rational prime \(q\) completely split in \(K\),
the ordinary locus of the universal abelian scheme is dense in each
geometric component of the special fiber of \(S\).
\end{lemma}

\begin{proof}
\emph{Models and Albanese generation.}
Spread the finitely many schemes, morphisms, and universal data over
an open part of the ring of integers of \(K\).  One can obtain the
proper model of \(S\) by taking its closure in a projective
compactification of the full parameter scheme used in the proof of
Lemma~\ref{lem:reference-moduli}.  Its intersection with the boundary
has proper closed image on the integer base.  That image misses the
generic point because \(S\) is proper and lies in the parameter
scheme.  Removing the finitely many primes in the image leaves the
closure inside the full parameter scheme.  Thus its geometric points
are still distinguished by their polarized, acted-on, marked tuples,
including the graph coordinates.  Further shrinking gives smooth
proper models with geometrically connected component fibers, good
abelian models, and the extended Albanese maps.

For a component \(S_j\), let
\(\Delta_j=a_j(S_j)-a_j(S_j)\subset A_j\).  It is closed,
irreducible, symmetric, and contains zero.  The successive sum
images \(n\Delta_j\) are nested irreducible closed sets.  Their
dimensions eventually stabilize; then equality of two successive
sets implies stability under addition by \(\Delta_j\).  Symmetry
shows that the stable image is the abelian subvariety generated by
\(\Delta_j\), which is \(A_j\) by the Albanese property.  Thus for
some fixed \(n\) the proper sum map
\[
(S_j\times S_j)^n\longrightarrow A_j
\]
is surjective.  Spread this map together with the other data.  Its
closed image contains the generic fiber of the flat abelian scheme;
that fiber is dense, so the image is the whole abelian scheme.
Consequently the sum map is surjective on every geometric fiber.

\emph{Ordinary points.}
For each characteristic-zero component, we now choose a point whose
reduction at the split primes retained below is ordinary.  In its
rational Hermitian three-space choose, by weak approximation, an \(E\)-line
negative at \(\mu\), and diagonalize its orthogonal complement.
At the period whose negative line is this coordinate line, the three
rational orthogonal projectors split the rational homology
into three \(E\)-lines and preserve the filtration.  Integral
multiples of them are endomorphisms of the corresponding abelian
variety.

These points and endomorphisms are defined over \(\Qbar\).  Indeed,
add their graphs as extra parameters with the Hilbert polynomials
at the selected point.  This is a finite-type parameter scheme over
a number field.  Locally the projector matrices are fixed by
rigidity.  A negative line preserved by all three projectors must
be the unique negative coordinate line, so the selected enhanced
point is isolated.  An isolated complex point of a finite-type
scheme over a number field has algebraic coordinates.  Adjoin the
finitely many fields of definition to \(K\), spread these points
and endomorphisms, and exclude their projector and isogeny
denominators.

Now let \(q\) be a remaining completely split prime.  At each
selected point the integral projectors and the idempotents of
\(\mathcal O_E\otimes\Z_q\) split the \(q\)-divisible group into
pieces of height one.  Their heights are the corresponding
characteristic-zero Tate ranks and stay constant under specialization.
A height-one \(q\)-divisible group over an algebraically closed
field has dimension zero or one.  In the first case it is étale;
in the second its dual has dimension zero, so it is multiplicative.
Thus each selected reduction is ordinary.  Ordinariness is open,
and a smooth geometrically connected fiber is geometrically
irreducible.  This nonempty open is dense on each component.
Chebotarev's theorem supplies infinitely many completely split
rational primes after the finite exclusions.
\end{proof}

Choose such a rational prime \(q\), also prime to the polarization
degree, \(N\), the discriminant of \(E\), and all exceptional local
data of the finite theta inputs.  Fix an embedding
\[
\iota:\Qbar\hookrightarrow\overline{\Q}_q
\]
and normalize \(\nu_\iota\) by \(\nu_\iota(q)=1\).  The resulting
completion of \(K\) is \(\Q_q\), with residue field \(\Fq\).
For \(\tau\in G_E\), let \(\mathfrak p_\tau\) be the prime of \(E\)
defined by \(\iota\circ\tau\).  Thus for \(x\in E^\times\),
\[
\nu_\iota(\tau(x))=\operatorname{val}_{\mathfrak p_\tau}(x).
\]
Select one prime from each conjugate pair, including
\begin{equation}\label{eq:active-prime}
w=\mathfrak p_c,\qquad \bar w=\mathfrak p_1.
\end{equation}
Write \(\mathcal J\) for this set of selected blocks.  The
\(w\)-block has height three and Lie rank one; every other selected
block has height three and Lie rank zero or three.  The opposite
blocks are their polarization duals.  These rank assertions persist
on the model because the corresponding idempotent Lie and torsion
summands are locally free.  Choosing \(w=\mathfrak p_c\), rather
than \(\mathfrak p_1\), will fix the conjugate-valuation signs.

\subsection{Algebraic Hecke operators on the Albanese}\label{subsec:algebraic-hecke}

Choose bases for the selected \(q\)-adic lattices and the opposite
dual lattices.  Since the polarization is prime to \(q\), the split
local similitude group and its integral subgroup are
\[
\mathscr G_q=
\left(\prod_{j\in\mathcal J}\GL_3(\Q_q)\right)\times\Q_q^\times,
\qquad
\mathscr K_q=
\left(\prod_{j\in\mathcal J}\GL_3(\Z_q)\right)\times\Z_q^\times.
\]
The last coordinate is the common polarization multiplier.
Use the raw spherical Hecke algebra with
\(\operatorname{vol}(\mathscr K_q)=1\).  Its integral basis consists
of double-coset characteristic functions; its rational extension
will act on the Albanese up to isogeny.  A double coset is specified
by a multiplier exponent \(h\in\Z\) and, on each selected block,
an unordered triple of elementary exponents
\((a_{j,1},a_{j,2},a_{j,3})\).
For a representative \(g_j\) with those exponents, our target
position is the overlattice \(g_j^{-1}\Lambda_j\); branches are
represented modulo the integral subgroup on the left.

First suppose that
\begin{equation}\label{eq:positive-hecke-range}
h\geq0,\qquad 0\leq a_{j,k}\leq h.
\end{equation}
On the \(j\)-block of \(B[q^h]\), choose a subgroup of the indicated
elementary divisors.  On the opposite block choose its annihilator
under the perfect \(q^h\)-torsion pairing.  The total subgroup \(H\)
is maximal isotropic: each conjugate pair contributes order
\(q^{3h}\), so its total order is \(q^{3sh}\).
Its quotient \(f:B\to B'=B/H\) is required
to have polarization
\begin{equation}\label{eq:hecke-polarization}
f^\vee\lambda'f=q^h\lambda.
\end{equation}
We verify that this is an algebraic construction in families,
including the polarization homomorphism.

Over the characteristic-zero parameter scheme the torsion and the
subgroups are finite étale.  The translation quotient exists as an
fppf scheme quotient for the free finite-locally-free action
\cite[Tags~07S6 and~07S7]{StacksProject}.  Its group law descends,
and the finite étale quotient of the abelian scheme is again an
abelian scheme.

On a complex fiber the polarization in
\eqref{eq:hecke-polarization} is seen directly on lattices.  If
\(\Lambda_j\subset\Lambda_j'\subset q^{-h}\Lambda_j\) is the
selected target overlattice and the multiplier is \(q^h\), its
opposite lattice is
\[
q^{-h}(\Lambda_j')^\vee.
\]
The \(q^h\)-scaled alternating form is integral and perfect on this
pair at \(q\), by the annihilator condition.  At every other prime
the \(q\)-isogeny changes no elementary-divisor valuation.  The
scaled form is positive, so it gives a Riemann polarization on the
quotient of the original type.  More generally, the field-level
isotropic descent criterion gives the unique symmetric polarization
homomorphism on every geometric fiber
\cite[Chapter~XI, Proposition~11.25]{EdixhovenVanDerGeerMoonen2011Polarizations}.

Here is the relative descent from this fiber calculation.  Since
\(H\subset B[q^h]\), multiplication by \(q^h\) descends to
\(g:B'\to B\), with
\[
gf=[q^h]_B,\qquad fg=[q^h]_{B'}.
\]
Put \(\beta=g^\vee\lambda g:B'\to B'^\vee\).
On each geometric fiber the preceding polarization satisfies
\(\beta=q^h\lambda'\), so \(\beta\) kills \(B'[q^h]\).
The correspondence
parameter scheme is smooth and reduced, and \(B'[q^h]\) is finite
étale over it, hence reduced.  Vanishing on its geometric points is
therefore vanishing of the relative morphism.  The fppf quotient by
\(B'[q^h]\) now factors \(\beta\) uniquely through \([q^h]_{B'}\).
It gives a relative homomorphism \(\lambda'\) with
\(\beta=q^h\lambda'\).  Its fibers are the symmetric polarizations
just constructed, and cancellation of isogenies gives
\eqref{eq:hecke-polarization}.  This also proves uniqueness.  We
have descended the polarization homomorphism, not a particular
line bundle.

The \(\mathcal O_E\)-action descends because \(H\) is stable.
Cancelling \(f\) in \eqref{eq:hecke-polarization} shows that its
Rosati adjoint remains \(c\).  The characteristic-zero Lie ranks
are unchanged, and \(f\) transports the full \(N\)-basis.  Thus
the quotient is a point of the same full tuple space.  The subgroup
choices form a finite étale scheme over \(S\); the quotient
construction is its algebraic target map.  Distinct kernels remain
distinct branches even when their target tuples are isomorphic.

\begin{lemma}\label{lem:hecke-albanese}
The rational spherical Hecke algebra of
\((\mathscr G_q,\mathscr K_q)\) acts in \(\End_K(A_S)\otimes\Q\) and,
by pullback, on \(H^1(S,\Q)\).  A basis operator in the positive
range \eqref{eq:positive-hecke-range} is an integral homomorphism
defined over \(K\).  Its pullback on one-forms is the sum of
pullbacks over its target branches.
\end{lemma}

\begin{proof}
Let \(C\to S\) be one of the finite étale source correspondences.
On a source component, compose each target branch with the
appropriate Albanese map and sum in \(A_S\).  This sum descends
from a cover trivializing the branches, because it is invariant
under their permutations.  It is a morphism \(n_j:S_j\to A_S\).
Subtracting \(n_j\) at the basepoint and applying the Albanese
universal property gives an integral homomorphism
\(A_j\to A_S\).  On differences it has the exact formula
\begin{equation}\label{eq:hecke-albanese-differences}
T\bigl(a_j(x)-a_j(y)\bigr)
 =\sum_{z\in C_x}a_{j(z)}(t(z))
  -\sum_{z\in C_y}a_{j(z)}(t(z)),
\end{equation}
where \(t\) is the target map, \(j(z)\) indexes the component
containing \(t(z)\), fibers are counted with their branches, and
\(a_{j(z)}\) is included in \(A_S\).  These
homomorphisms on the factors define the operator on \(A_S\).
Their pullback on one-forms is the branch sum, since a constant
translation contributes no differential.

Composition has the raw spherical convolution multiplicities.
For a selected lattice \(\Lambda_j\), a branch is an overlattice
\(g_j^{-1}\Lambda_j\) with a representative modulo the integral
subgroup on the left.  If its cumulative multiplier is \(q^h\),
the opposite target lattice is forced to be
\(q^{-h}(\Lambda_j')^\vee\).  Consequently a chain of quotient
choices is precisely a chain of selected overlattices, as in
spherical convolution.  The opposite lattices, polarizations, and
transported markings agree for the same composite.  The numbers
of branches landing in each target component are locally constant
on a connected source.  Thus the component basepoint translations
cancel when \eqref{eq:hecke-albanese-differences} is composed.

Consider the central position with all selected exponents equal to
one and multiplier exponent two.  It is quotient by \(B[q]\);
under the multiplication-by-\(q\) identification its target has
the same polarization and action and marking \(q\eta_N\).  It is
therefore a finite-order marking automorphism, since \(q\) is
invertible modulo \(N\).  Adding a sufficiently large power of
this position to any integer exponent position puts it in
\eqref{eq:positive-hecke-range}.  Define its operator by clearing
that power with the inverse marking automorphism.  The positive
convolution identity shows that this definition is independent
of the clearing power.  This gives the full rational spherical
action.  The spherical algebra is commutative, as follows also
from the Satake isomorphism \cite[\S3]{Gross1998}.
\end{proof}
